% =====================================================================
%  Cac hinh bo sung cho Chuong 2 va Chuong 3
%  Bien dich:  xelatex hinh_ch23.tex
% =====================================================================
\documentclass[11pt]{article}
\usepackage[margin=1.6cm,paperwidth=19cm,paperheight=30cm]{geometry}
\usepackage{fontspec}
\setmainfont{Liberation Serif}   % Overleaf: doi thanh "Times New Roman"
\usepackage{amsmath}
\usepackage{tikz}
\usepackage{pgfplots}
\pgfplotsset{compat=1.18}
\usetikzlibrary{positioning,fit,shapes.geometric,shapes.symbols,backgrounds,arrows.meta,calc,decorations.pathreplacing}
\pagestyle{empty}
\setlength{\parindent}{0pt}

\tikzset{
  box/.style ={draw,fill=gray!12,align=center,inner sep=3pt,line width=0.5pt},
  sbox/.style={draw,align=center,inner sep=3pt,font=\scriptsize,line width=0.5pt},
}
\newcommand{\chuthich}[2]{\vspace{1.5mm}{\footnotesize #1\\ \textit{Nguồn: #2}}\vspace{7mm}}

\begin{document}
\centering

% =====================================================================
% HINH A -- muc 2.8.1 : tran toc do do nhieu luong rieng
% =====================================================================
\begin{tikzpicture}[font=\footnotesize]
\begin{axis}[
  name=left, width=8.2cm, height=6.0cm,
  xlabel={Tỷ số tín hiệu trên tạp âm $g_kP_{\mathrm{priv}}/\sigma^2$ (dB)},
  ylabel={$R_{p,k}$ (bit/s/Hz)},
  xmin=0,xmax=30, ymin=0, ymax=10.5, ytick={0,2,4,6,8,10},
  grid=both, grid style={gray!25}, tick label style={font=\scriptsize},
  label style={font=\scriptsize},
  legend style={font=\scriptsize,at={(0.03,0.97)},anchor=north west,draw=gray!50},
  legend cell align=left,
]
% M = 1 : toan bo cong suat rieng don vao mot luong
\addplot[very thick,blue!70!black,domain=0:30,samples=120]
   {log2(1 + 10^(x/10))};
\addlegendentry{$M=1$ (tập trung)}
% M = 2
\addplot[very thick,orange!85!black,dashed,domain=0:30,samples=120]
   {log2(1 + 10^(x/10)/(1*10^(x/10) + 2))};
\addlegendentry{$M=2$ (chia đều)}
% M = 4
\addplot[very thick,violet,dotted,domain=0:30,samples=120]
   {log2(1 + 10^(x/10)/(3*10^(x/10) + 4))};
\addlegendentry{$M=4$ (chia đều)}
% duong tiem can
\addplot[gray,thin,domain=0:30,samples=2]{1.0};
\addplot[gray,thin,domain=0:30,samples=2]{0.415};
\node[font=\tiny,anchor=west,gray!45!black,align=left] at (axis cs:13.5,1.9)
     {đường xám: trần tiệm cận $\log_2\!\big(M/(M-1)\big)$};
\end{axis}

\begin{axis}[
  at={(left.east)}, anchor=west, xshift=1.1cm,
  width=6.6cm, height=6.0cm,
  xlabel={Số luồng riêng hoạt động $M$},
  ylabel={$R_{p,k}$ tiệm cận (bit/s/Hz)},
  ybar, bar width=9pt, xmin=1.4,xmax=8.6, ymin=0,ymax=1.2,
  xtick={2,3,4,5,6,7,8},
  grid=both, grid style={gray!25}, tick label style={font=\scriptsize},
  label style={font=\scriptsize},
  nodes near coords, every node near coord/.append style={font=\tiny,/pgf/number format/.cd,fixed,precision=3,use comma},
]
\addplot[fill=blue!35,draw=blue!60!black] coordinates
  {(2,1.000)(3,0.585)(4,0.415)(5,0.322)(6,0.263)(7,0.222)(8,0.193)};
\end{axis}
\end{tikzpicture}

\chuthich{Hình 2.4: Trần tốc độ luồng riêng do nhiễu giữa các luồng riêng trong hệ SISO.
Trái: $R_{p,k}$ theo tỷ số tín hiệu trên tạp âm. Phải: giá trị tiệm cận khi $M$ luồng riêng chia đều công suất.}
{Tác giả xây dựng từ công thức (2.29)--(2.32) (2026)}

% =====================================================================
% HINH B -- muc 2.8.4 : vai tro cua can pha
% =====================================================================
\begin{tikzpicture}[font=\footnotesize,>=Stealth,scale=1.0]

% ---------- panel trai: pha lech ----------
\begin{scope}[shift={(0,0)}]
  \draw[gray!45,->] (-0.6,0)--(3.6,0);
  \draw[gray!45,->] (0,-1.9)--(0,1.9);
  \draw[blue!65,line width=1pt,->] (0,0) -- ++(20:0.75);
  \draw[blue!65,line width=1pt,->] (20:0.75) -- ++(100:0.75);
  \draw[blue!65,line width=1pt,->] ($(20:0.75)+(100:0.75)$) -- ++(175:0.75);
  \draw[blue!65,line width=1pt,->] ($(20:0.75)+(100:0.75)+(175:0.75)$) -- ++(250:0.75);
  \draw[blue!65,line width=1pt,->] ($(20:0.75)+(100:0.75)+(175:0.75)+(250:0.75)$) -- ++(320:0.75);
  \coordinate (T1) at ($(20:0.75)+(100:0.75)+(175:0.75)+(250:0.75)+(320:0.75)$);
  \draw[red!75!black,line width=1.5pt,->] (0,0) -- (T1);
  \node[font=\scriptsize,red!75!black,anchor=west] at (0.30,-0.30) {$\big|\sum_n\big|$ nhỏ};
  \node[font=\scriptsize,anchor=north] at (1.3,-2.15) {(a) Pha chưa căn chỉnh};
\end{scope}

% ---------- panel phai: da can pha ----------
\begin{scope}[shift={(6.6,0)}]
  \draw[gray!45,->] (-0.6,0)--(4.6,0);
  \draw[gray!45,->] (0,-1.9)--(0,1.9);
  \foreach \i in {0,1,2,3,4}{
     \draw[blue!65,line width=1pt,->] (\i*0.75,0.30) -- (\i*0.75+0.75,0.30);
  }
  \draw[red!75!black,line width=1.5pt,->] (0,-0.30) -- (3.75,-0.30);
  \node[font=\scriptsize,red!75!black,anchor=north] at (1.9,-0.38) {$\big|\sum_n\big|$ lớn};
  \node[font=\scriptsize,anchor=north] at (1.9,-2.15) {(b) Sau khi căn pha $\bar\theta_n=-\angle z_n$};
\end{scope}

\node[font=\scriptsize,align=center,anchor=north] at (4.2,-2.75)
  {Mỗi mũi tên là một thành phần ghép tầng $h^{*}_{r,k,n}g_n e^{j\theta_n}$ qua một phần tử STAR--RIS};
\end{tikzpicture}

\chuthich{Hình 2.5: Ý nghĩa của căn pha đối với độ lợi kênh hiệu dụng}
{Tác giả xây dựng dựa trên [6], [7] (2026)}

\clearpage

% =====================================================================
% HINH C -- muc 3.2 : cau truc vec-to trang thai
% =====================================================================
\begin{tikzpicture}[font=\footnotesize,>=Stealth]
\def\h{0.85}
% cac khoi
\node[sbox,fill=blue!12,  minimum width=2.1cm,minimum height=\h cm,anchor=west] (b1) at (0,0)    {$\mathrm{Re}/\mathrm{Im}\ h_{d,k}$};
\node[sbox,fill=green!14, minimum width=2.6cm,minimum height=\h cm,anchor=west] (b2) at (2.1,0)  {$\mathrm{Re}/\mathrm{Im}\ g_n$};
\node[sbox,fill=orange!16,minimum width=5.2cm,minimum height=\h cm,anchor=west] (b3) at (4.7,0)  {$\mathrm{Re}/\mathrm{Im}\ h_{r,k,n}$};
\node[sbox,fill=violet!14,minimum width=2.3cm,minimum height=\h cm,anchor=west] (b4) at (9.9,0)  {chỉ báo phía $s_k$};

% so chieu duoi tung khoi
\node[font=\scriptsize,anchor=north] at (1.05,-0.50) {$2K$};
\node[font=\scriptsize,anchor=north] at (3.40,-0.50) {$2N$};
\node[font=\scriptsize,anchor=north] at (7.30,-0.50) {$2KN$};
\node[font=\scriptsize,anchor=north] at (11.05,-0.50) {$K$};
\node[font=\tiny,anchor=north,gray!60!black] at (1.05,-0.88) {$=8$};
\node[font=\tiny,anchor=north,gray!60!black] at (3.40,-0.88) {$=64$};
\node[font=\tiny,anchor=north,gray!60!black] at (7.30,-0.88) {$=256$};
\node[font=\tiny,anchor=north,gray!60!black] at (11.05,-0.88) {$=4$};

% he so chuan hoa phia tren
\node[font=\scriptsize,anchor=south] at (1.05,0.52) {$\div\,0{,}35$};
\node[font=\scriptsize,anchor=south] at (3.40,0.52) {$\div\,1{,}0$};
\node[font=\scriptsize,anchor=south] at (7.30,0.52) {$\div\,0{,}75$};
\node[font=\scriptsize,anchor=south] at (11.05,0.52) {$\{0,1\}\!\to\!\{-1,1\}$};

% ngoac tong
\draw[decorate,decoration={brace,amplitude=5pt,mirror},gray!70]
      (0,-1.35) -- (12.2,-1.35);
\node[font=\scriptsize,anchor=north] at (6.1,-1.80)
     {$d_s = 2(K+N+KN)+K = 332$ \quad với $K=4,\ N=32$};

% buoc cat bien
\node[box,anchor=west,font=\scriptsize] (clip) at (0,-2.75) {Cắt biên về $[-8,\,8]$};
\node[box,anchor=west,font=\scriptsize] (net)  at (4.2,-2.75) {Mạng nơ-ron (Actor / Critic)};
\draw[->] ($(clip.north)+(0,0.60)$) -- (clip.north);
\draw[->] (clip.east) -- (net.west);
\end{tikzpicture}

\chuthich{Hình 3.2: Cấu trúc và chuẩn hóa theo khối của véc-tơ trạng thái}
{Tác giả xây dựng theo công thức (3.1), (3.2) và mã nguồn (2026)}

% =====================================================================
% HINH D -- ve lai Hinh 3.2 : bo giai ma hanh dong
% =====================================================================
\begin{tikzpicture}[font=\footnotesize,>=Stealth]

% ---------- vec-to hanh dong, xep doc ----------
\node[font=\scriptsize,anchor=south] at (0.6,0.45) {$a\in[-1,1]^{d_a}$};
\foreach \y/\lbl/\dim in {0/{$a_c$}/1, -1.1/{$a_p$}/K, -2.2/{$a_\eta$}/K,
                          -3.3/{$a_\beta$}/N, -4.40/{$a_{\theta^T}$}/N, -5.05/{$a_{\theta^R}$}/N}{
   \node[sbox,fill=gray!18,minimum width=1.2cm,minimum height=0.55cm,anchor=west] at (0,\y) {\lbl};
   \node[font=\tiny,gray!55!black,anchor=east] at (-0.15,\y) {\dim};
}
\node[font=\scriptsize,anchor=north] at (0.6,-5.55) {$d_a=2K+1+3N$};

% ---------- khoi giai ma ----------
\node[box,anchor=west,text width=4.9cm,font=\scriptsize,minimum height=0.8cm] (d1) at (2.4,0)
     {Ánh xạ tuyến tính\\ $\rho_c=0{,}70+0{,}29\,\tfrac{a_c+1}{2}$};
\node[box,anchor=west,text width=4.9cm,font=\scriptsize,minimum height=0.8cm] (d2) at (2.4,-1.1)
     {Softmax độ sắc $\kappa=10$};
\node[box,anchor=west,text width=4.9cm,font=\scriptsize,minimum height=0.8cm] (d3) at (2.4,-2.2)
     {Chiếu Euclid lên đơn hình $\Delta_K$};
\node[box,anchor=west,text width=4.9cm,font=\scriptsize,minimum height=0.8cm] (d4) at (2.4,-3.3)
     {Ánh xạ tuyến tính $[-1,1]\to[0,1]$};
\node[box,anchor=west,text width=4.9cm,font=\scriptsize,minimum height=0.8cm] (d5) at (2.4,-4.725)
     {Pha tham chiếu $\bar\theta_n=-\angle z_n$\\ cộng phần hiệu chỉnh học được};

% ---------- bien vat ly ----------
\node[sbox,anchor=west,fill=blue!10,text width=5.3cm,font=\scriptsize,minimum height=0.8cm] (v1) at (8.1,0)
     {$p_c=\rho_c P_{\max}$, \ $\rho_c\in[0{,}70;\,0{,}99]$};
\node[sbox,anchor=west,fill=blue!10,text width=5.3cm,font=\scriptsize,minimum height=0.8cm] (v2) at (8.1,-1.1)
     {$p_k=(1-\rho_c)P_{\max}w_k$ \ (gần one-hot)};
\node[sbox,anchor=west,fill=blue!10,text width=5.3cm,font=\scriptsize,minimum height=0.8cm] (v3) at (8.1,-2.2)
     {$\eta_k\ge 0$, \ $\textstyle\sum_k \eta_k = 1$};
\node[sbox,anchor=west,fill=blue!10,text width=5.3cm,font=\scriptsize,minimum height=0.8cm] (v4) at (8.1,-3.3)
     {$\beta^T_n\in[0,1]$, \ $\beta^R_n=1-\beta^T_n$};
\node[sbox,anchor=west,fill=blue!10,text width=5.3cm,font=\scriptsize,minimum height=0.8cm] (v5) at (8.1,-4.725)
     {$\theta^T_n,\ \theta^R_n=\bar\theta_n+0{,}25\pi\,a_{\theta}$};

% ---------- noi day: hang ngang, khong cat nhau ----------
\foreach \i/\y in {1/0, 2/-1.1, 3/-2.2, 4/-3.3}{
   \draw[->] (1.2,\y) -- (d\i.west);
}
\draw (1.2,-4.40) -- (1.85,-4.40);
\draw (1.2,-5.05) -- (1.85,-5.05);
\draw (1.85,-4.40) -- (1.85,-5.05);
\draw[->] (1.85,-4.725) -- (d5.west);

\foreach \i in {1,...,5}{ \draw[->] (d\i.east) -- (v\i.west); }

\node[font=\scriptsize,anchor=north west,text width=13.4cm] at (0,-6.05)
  {Mọi ràng buộc công suất, đơn hình và bảo toàn năng lượng được bảo đảm ngay tại bộ giải mã,
   nên không cần phép chiếu hay hạng phạt bổ sung sau khi hành động đã được sinh.};
\end{tikzpicture}

\chuthich{Hình 3.3: Bộ giải mã hành động có cấu trúc vật lý}
{Tác giả xây dựng theo \texttt{action.py} (2026)}

\clearpage

% =====================================================================
% HINH E -- muc 3.4.2 : anh huong cua do sac softmax
% =====================================================================
\begin{tikzpicture}[font=\footnotesize]
\begin{axis}[
  width=13.5cm, height=6.4cm,
  ybar=2pt, bar width=11pt,
  xlabel={Chỉ số UE $k$}, ylabel={Trọng số công suất riêng $w_k$},
  symbolic x coords={UE1,UE2,UE3,UE4},
  xtick=data, ymin=0, ymax=1.08,
  grid=both, grid style={gray!25},
  tick label style={font=\scriptsize}, label style={font=\scriptsize},
  legend style={font=\scriptsize,at={(0.98,0.97)},anchor=north east,draw=gray!50},
  legend cell align=left,
  nodes near coords, every node near coord/.append style=
     {font=\tiny,/pgf/number format/.cd,fixed,precision=3,use comma},
]
\addplot[fill=gray!35,draw=gray!60!black]      coordinates {(UE1,0.358)(UE2,0.265)(UE3,0.217)(UE4,0.161)};
\addlegendentry{$\kappa=1$}
\addplot[fill=orange!45,draw=orange!70!black]  coordinates {(UE1,0.756)(UE2,0.169)(UE3,0.062)(UE4,0.014)};
\addlegendentry{$\kappa=5$}
\addplot[fill=blue!45,draw=blue!70!black]      coordinates {(UE1,0.946)(UE2,0.047)(UE3,0.006)(UE4,0.000)};
\addlegendentry{$\kappa=10$ (luận văn)}
\end{axis}
\node[font=\scriptsize,anchor=north] at (6.8,-1.35)
     {Cùng một đầu ra Actor $a_p=[0{,}5;\ 0{,}2;\ 0{,}0;\ -0{,}3]$};
\end{tikzpicture}

\chuthich{Hình 3.4: Ảnh hưởng của hệ số độ sắc $\kappa$ đến phân bổ công suất luồng riêng}
{Tác giả xây dựng theo công thức (3.6) (2026)}

\end{document}
